% BHU PYQ -- Professional Exam Template
\documentclass[11pt,a4paper]{article}

\usepackage[a4paper,top=2.5cm,bottom=2.5cm,left=2.8cm,right=2.8cm,headheight=14pt]{geometry}
\usepackage{amsmath,amssymb}
\usepackage[T1]{fontenc}
\usepackage[utf8]{inputenc}
\usepackage{lmodern}
\usepackage{microtype}
\usepackage{fancyhdr}
\usepackage{enumitem}
\usepackage{hyperref}
\usepackage{lastpage}
\usepackage{parskip}

\hypersetup{colorlinks=true,urlcolor=black,linkcolor=black,
  pdftitle={B.Sc. (Hons.) Botany Sem-VI BOB-602 -- BHU PYQ}}

\pagestyle{fancy}
\fancyhf{}
\renewcommand{\headrulewidth}{0.4pt}
\renewcommand{\footrulewidth}{0.4pt}
\fancyhead[L]{\small   \textit{Banaras Hindu University}}
\fancyhead[R]{\small   \textit{Botany Paper: BOB-602}}
\fancyfoot[C]{\small Page  \thepage\ of \pageref{LastPage}}


\newlist{parts}{enumerate}{2}
\setlist[parts,1]{label=(\alph*),leftmargin=*,itemsep=4pt,topsep=3pt}
\setlist[parts,2]{label=(\roman*),leftmargin=*,itemsep=2pt,topsep=2pt}

\newcommand{\pts}[1]{\hfill   \textit{[#1]}}


\begin{document}


\begin{center}
  {\large\bfseries BANARAS HINDU UNIVERSITY}\\[2pt]
  {
\normalsize B.Sc. (Hons.) Semester-VI Examination, 2022-23}\\[6pt]
  \rule{0.8\linewidth}{0.5pt}\\[6pt]
  {\large\bfseries Botany}\\[4pt]
  {\large\bfseries Paper: BOB-602 --- Microbiology and Plant Pathology}\\[4pt]
  {
\normalsize Time: Three Hours \hfill Maximum Marks: 70}
\end{center}

\rule{\linewidth}{0.4pt}

\smallskip



\noindent   \textit{Instructions: Answer five questions in all, selecting two questions from each Section-A and B. Question No. 1 is compulsory. The figures in the right-hand margin indicate marks.}

\bigskip

\medskip


\noindent   \textbf{Question 1.} Answer the following in not more than 50 words each:\pts{$1  \times 10 = 10$}

\begin{parts}
  \item What is Hfr cell?
  \item What is   \textit{Halobacterium}?
  \item What is the difference between lag and log phases of bacterial growth?
  \item Mention the contribution of Carl Woese.
  \item What is Mycoplasma?
  \item What are biotrophic pathogens? Give an example.
  \item Mention the contribution of Henry Phipps in Indian Plant Pathology.
  \item What is PTI?
  \item Name two insect fungi that are used as bio-control agents.
  \item Name a plant disease caused by phytoplasma.
\end{parts}

\bigskip

\begin{center}
   \textbf{SECTION A}
\end{center}
\medskip

\medskip


\noindent   \textbf{Question 2.} Describe the mechanism of transduction in bacteria with well-labelled diagrams.\pts{15}

\medskip


\noindent   \textbf{Question 3.} Write short notes on any two of the following:\pts{$7.5  \times 2 = 15$}

\begin{parts}
  \item Free-living and symbiotic microorganisms
  \item Cyanophage
  \item Lactic acid fermentation by bacteria
\end{parts}

\medskip


\noindent   \textbf{Question 4.} Give a detailed account of the interaction of   \textit{Rhizobium} with a leguminous plant.\pts{15}

\medskip


\noindent   \textbf{Question 5.} Write short notes on any two of the following:\pts{$7.5  \times 2 = 15$}

\begin{parts}
  \item Compare between Archaea and Bacteria
  \item Compare between Lytic and Lysogenic cycle
  \item Plant growth promoting rhizobacteria
\end{parts}

\bigskip

\begin{center}
   \textbf{SECTION B}
\end{center}
\medskip

\medskip


\noindent   \textbf{Question 6.} Briefly describe the various modes of infection in plants caused by fungal pathogens with suitable examples.\pts{15}

\medskip


\noindent   \textbf{Question 7.} What are the different types of induced biochemical defense in plants? Describe briefly the role of PR proteins with suitable examples.\pts{15}

\medskip


\noindent   \textbf{Question 8.} Answer any two of the following in brief:\pts{$7.5  \times 2 = 15$}

\begin{parts}
  \item Causal organism and symptoms of the wilt of tomato
  \item Modes of the infection of fungal pathogens
  \item Downy mildew of crucifers
\end{parts}

\medskip


\noindent   \textbf{Question 9.} Answer any two of the following in brief:\pts{$7.5  \times 2 = 15$}

\begin{parts}
  \item Smut of Bajra
  \item Integrated disease management
  \item Transmission pathways of the sugarcane mosaic disease
\end{parts}

\vfill
\rule{\linewidth}{0.4pt}

\begin{center}\small
\href{https://bhu-exam.vercel.app/}{Click Here}\\[4pt]\href{https://me-aryan.vercel.app/}{Click Here}\\[4pt]\href{https://quashan-me.vercel.app/}{Click Here}
\end{center}

\end{document}
